\documentclass[12pt,a4paper]{article}
\usepackage[margin=2.5cm]{geometry}
\usepackage{fontenc}
\usepackage{amsmath}
\usepackage{listings}
\usepackage{xcolor}
\usepackage{hyperref}
\usepackage{tabularx}
\usepackage{longtable}
\usepackage{booktabs}
\hypersetup{colorlinks=true, linkcolor=blue, urlcolor=blue}
\lstset{
    basicstyle=\ttfamily\small,
    keywordstyle=\color{blue},
    commentstyle=\color{gray},
    breaklines=true,
    language=C
}
\title{SLFD Specification\\\large{Simple Library Functions Define}}
\author{Version 10.2}
\date{\today}

\begin{document}

\maketitle

% Table of contents disabled to avoid TOC write issues during build
%\tableofcontents
%\newpage

\section{Knowledge}
Welcome to Simple Library Functions Define (SLFD). 
This is an Interface Specification of E-comOS or other OSes that implement \textbf{SLFD}.
Please note we support C language only for now, but we may support other languages in the future (not include C++, never).
If you are a C++ coder, and you want to use SLFD, please use 
\begin{lstlisting}
    extern "C" {
        /* SLFD code here */
    }
\end{lstlisting}
to do your works.

\section{Scope}
This document defines the core operating system interface known as \textbf{SLFD} (Simple Library Functions Define). It provides system calls for process management, memory allocation, file/device I/O, pipes, network communications, event multiplexing (hear interface), inter‑process communication (IPC), time operations, environment variables, file system operations, command execution, thread management, zero‑copy I/O, asynchronous I/O, and additional features such as user/group management, shared memory, resource limits, and privilege elevation.

The interface is designed to be simple and extensible. Security policies (e.g., privilege levels, access control) are described in Appendix A (ULA - User Level Assessment) and are fully integrated into the core API using integer bitmasks.

All file system paths in this specification follow the SLFD FD layout defined in Appendix B, which includes top‑level directories \texttt{/bin}, \texttt{/dev}, \texttt{/conf}, \texttt{/usr}, \texttt{/lib}, \texttt{/tmp}, \texttt{/sys}, and \texttt{/mnt}.

Kernel Architecture Note: This specification does not mandate a particular kernel architecture. Both monolithic kernels and microkernels can implement SLFD. The behaviour described for each function is mandatory; the underlying implementation may use direct trapping, user‑space servers, or any combination thereof.

\subsection{Design Philosophy}
SLFD aims to provide a consistent, minimal, and secure interface. Key principles:
\begin{itemize}
    \item \textbf{Separation of creation and execution}: Process objects are created with \texttt{proc\_new} and later started with \texttt{proc\_start}, allowing configuration in between.
    \item \textbf{Explicit resource declaration}: Processes must declare their needed capabilities via \texttt{proc\_setflags} before starting (cannot be elevated later without \texttt{hlud}).
    \item \textbf{Unified I/O for local resources}: Files, devices, and pipes use \texttt{open}/\texttt{read}/\texttt{write}/\texttt{close}.
    \item \textbf{Dedicated network functions}: Network I/O uses separate calls (\texttt{net\_connect}, etc.) to avoid confusion with local file semantics.
    \item \textbf{Event multiplexing as a first-class citizen}: The \texttt{hear} interface allows waiting on any event source (I/O, IPC, signals, process state).
    \item \textbf{ULA privilege model}: Four levels encoded as bitmasks, replacing traditional UID/GID/mode\_t for better security and simplicity.
\end{itemize}

\section{Head Files}
Implementations shall provide the following header files. All headers are located under \texttt{/usr/include}.

\subsection[<slfd.h>]{\texorpdfstring{\texttt{<slfd.h>}}{<slfd.h>}}
Primary header, includes all other SLFD headers. Defines basic types, resource flags, signal numbers, I/O constants, and function prototypes.

\subsection[<slfd/ula.h>]{\texorpdfstring{\texttt{<slfd/ula.h>}}{<slfd/ula.h>}}
Defines ULA privilege levels and bitmask layout for file attributes.
\begin{lstlisting}
#define LV_HIGHEST   0
#define LV_GENERAL   1
#define LV_LOW       2
#define LV_GETTING   3

#define ULA_CREATOR_SHIFT     24
#define ULA_CALLER_SHIFT      16
#define ULA_FLAG_EXEC         (1<<8)
#define ULA_FLAG_HIDDEN       (1<<9)
#define ULA_FLAG_ARCHIVE      (1<<10)
#define ULA_TYPE_FILE         (0<<4)
#define ULA_TYPE_DIR          (1<<4)
#define ULA_TYPE_LINK         (2<<4)
#define ULA_TYPE_DEV          (3<<4)

#define ULA_MAKE_ATTR(creator, caller, flags, type) \
    (((creator) << ULA_CREATOR_SHIFT) | ((caller) << ULA_CALLER_SHIFT) | (flags) | (type))
\end{lstlisting}

\subsection[<slfd/fs.h>]{\texorpdfstring{\texttt{<slfd/fs.h>}}{<slfd/fs.h>}}
Extended file system operations (see Section~\ref{sec:fs_ext}).

\subsection[<slfd/thread.h>]{\texorpdfstring{\texttt{<slfd/thread.h>}}{<slfd/thread.h>}}
Threading primitives (see Section~\ref{sec:threads}).

\subsection[<slfd/performance.h>]{\texorpdfstring{\texttt{<slfd/performance.h>}}{<slfd/performance.h>}}
Zero‑copy and asynchronous I/O (implementation‑defined).

\subsection[<slfd/ipc_perf.h>]{\texorpdfstring{\texttt{<slfd/ipc\_perf.h>}}{<slfd/ipc_perf.h>}}
IPC performance tuning (implementation‑defined).

\subsection[<slfd/hear.h>]{\texorpdfstring{\texttt{<slfd/hear.h>}}{<slfd/hear.h>}}
Event multiplexing (hear interface), see Section~\ref{sec:hear}.


\section{Terminology and Conventions}

\subsection{Types}
\begin{lstlisting}
typedef int fid_t;              /* file/device/pipe/network identifier */
typedef int pid_t;              /* process identifier */
typedef long off_t;             /* file offset (signed 64-bit) */
typedef long ssize_t;           /* signed size (64-bit) */
typedef unsigned long size_t;   /* unsigned size (64-bit) */
typedef uint32_t thread_id;     /* thread identifier for IPC and threading */

typedef struct {
    int argc;
    char **argv;
    char **envp;
} parg_t;

struct stat {
    size_t st_size;
    uint32_t st_attr_mask;
    /* other fields implementation-defined */
};

struct utsname {
    char sysname[32];
    char nodename[32];
    char release[32];
    char version[32];
    char machine[32];
};

struct rlimit {
    unsigned long rlim_cur;
    unsigned long rlim_max;
};

#define IPC_MAX_DATA_SIZE 4096
typedef struct ipc_message {
    uint32_t type;
    uint32_t source;
    uint32_t target;
    uint32_t timestamp;
    uint32_t size;
    uint32_t sequence;
    uint8_t  data[IPC_MAX_DATA_SIZE];
} ipc_message_t;
\end{lstlisting}

\subsection{Standard Constants}
\begin{lstlisting}
/* Resource flags (bitmask) */
#define PROC_IO_READ    (1<<0)
#define PROC_IO_WRITE   (1<<1)
#define PROC_NETWORK    (1<<2)
#define PROC_MEMORY     (1<<3)

/* Signal numbers (PROCSIG_* are aliases to traditional signal numbers) */
#define PROCSIG_HUP      1
#define PROCSIG_INT      2
#define PROCSIG_QUIT     3
#define PROCSIG_ILL      4
#define PROCSIG_TRAP     5
#define PROCSIG_ABRT     6
#define PROCSIG_BUS      7
#define PROCSIG_FPE      8
#define PROCSIG_KILL     9
#define PROCSIG_USR1     10
#define PROCSIG_SEGV     11
#define PROCSIG_USR2     12
#define PROCSIG_PIPE     13
#define PROCSIG_ALRM     14
#define PROCSIG_TERM     15
#define PROCSIG_CHLD     17
#define PROCSIG_STOP     18
#define PROCSIG_CONT     19
#define PROCSIG_TSTP     20

/* io_seek whence */
#define SEEK_SET 0
#define SEEK_CUR 1
#define SEEK_END 2

/* mount flags */
#define MS_RDONLY     1
#define MS_NOSUID     2
#define MS_NODEV      4
#define MS_NOEXEC     8
#define MS_SYNCHRONOUS 16
#define MS_REMOUNT    32
#define MS_MANDLOCK   64
#define MS_DIRSYNC    128
#define MS_NOATIME    256
#define MS_NODIRATIME 512
#define MS_BIND       1024
#define MS_MOVE       2048
#define MS_REC        4096
#define MS_SILENT     8192
#define MS_POSIXACL   16384
#define MS_UNBINDABLE 65536
#define MS_PRIVATE    131072
#define MS_SLAVE      262144
#define MS_SHARED     524288
#define MS_RELATIME   1048576
#define MS_KERNMOUNT  2097152
#define MS_I_VERSION  4194304
#define MS_STRICTATIME 8388608
#define MS_LAZYTIME   16777216

/* umask constants */
#define UMASK_DEFAULT 0022
\end{lstlisting}

\subsection{Error Handling}
Functions return an integer (\texttt{int} or \texttt{ssize\_t}). Negative values indicate an error; the absolute value is one of the error codes defined in Section~\ref{sec:errors}. Non‑negative indicates success (or positive data like PID, \texttt{fid\_t}, bytes transferred). On failure, the global \texttt{errno} is also set to the absolute value of the error code, allowing use of standard error-checking idioms. Callers may check either the return value or \texttt{errno}.

\subsection{Conventions}
\begin{itemize}
    \item Strings are null‑terminated.
    \item Passing \texttt{NULL} where not allowed yields \texttt{-EINVAL}.
    \item File descriptors 0,1,2 are not pre‑opened unless \texttt{<slfd/openstd.h>} is used.
    \item File attributes are integer bitmasks using ULA macros.
\end{itemize}
\section{Process Management}

\subsection{proc\_new}
\begin{lstlisting}
pid_t proc_new(const char *args, uint32_t privilege_mask);
\end{lstlisting}

\subsubsection{Description}
Creates a new process object and allocates a unique PID. The process does not start executing immediately; a subsequent call to \texttt{proc\_start} is always required, regardless of whether \texttt{args} is \texttt{"n"} (new) or \texttt{"s"} (clone). This ensures a unified creation-startup flow.

\subsubsection{Parameters}
\begin{tabularx}{\textwidth}{|l|l|X|}
\hline
\textbf{Parameter} & \textbf{Type} & \textbf{Description} \\
\hline
\verb|args| & \verb|const char *| & Mode: \texttt{"n"} (new, no memory image) or \texttt{"s"} (clone, inherits a copy of parent's memory image). \\
\verb|privilege_mask| & \verb|uint32_t| & ULA privilege mask (see \texttt{<slfd/ula.h>}). \\
\hline
\end{tabularx}

\subsubsection{Return Value}
\begin{itemize}
    \item \textbf{Success}: Returns a positive PID.
    \item \textbf{Failure}: Returns a negative error code:
        \begin{itemize}
            \item \texttt{-EINVAL}: Invalid \texttt{args} string.
            \item \texttt{-ENOMEM}: Insufficient kernel memory for new process.
            \item \texttt{-EPERM}: Caller lacks privilege to create process with the given \texttt{privilege\_mask}.
        \end{itemize}
\end{itemize}

\subsubsection{Design Rationale}
Traditional UNIX \texttt{fork()} clones and runs immediately. SLFD separates creation from starting to allow configuration (e.g., via \texttt{proc\_setflags}) between them. Both \texttt{"n"} and \texttt{"s"} modes now require \texttt{proc\_start} to run, making the flow uniform and predictable.

\subsubsection{Notes for Callers}
\begin{itemize}
    \item The \texttt{privilege\_mask} must be constructed using \texttt{ULA\_MAKE\_ATTR} or similar macros.
    \item After \texttt{proc\_new}, you may call \texttt{proc\_setflags} to adjust resource flags before starting.
    \item The child process (for \texttt{"s"}) inherits the parent's memory image at the time of \texttt{proc\_new}, but does not execute until \texttt{proc\_start}.
\end{itemize}

\subsubsection{Changelog}
\begin{tabularx}{\textwidth}{|l|X|}
\hline
\textbf{Version} & \textbf{Change} \\
\hline
1.0 & Initial design: \texttt{proc\_new} only supported \texttt{"n"} mode. \\
3.0 & Added \texttt{"s"} mode for clone (fork) semantics, which ran immediately. \\
10.1 & Removed \texttt{flags} parameter; \texttt{"s"} mode now also requires \texttt{proc\_start}. Added \texttt{proc\_setflags} for separate flag configuration. \\
\hline
\end{tabularx}

\subsection{proc\_setflags}
\begin{lstlisting}
int proc_setflags(pid_t pid, uint32_t flags);
\end{lstlisting}

\subsubsection{Description}
Sets or modifies the resource flags for a process created with \texttt{proc\_new}. This must be called before \texttt{proc\_start}. After the process has started, flags cannot be changed.

\subsubsection{Parameters}
\begin{tabularx}{\textwidth}{|l|l|X|}
\hline
\verb|pid| & \verb|pid_t| & PID of the process. \\
\verb|flags| & \verb|uint32_t| & Bitmask of resource flags (e.g., \texttt{PROC\_IO\_READ}). \\
\hline
\end{tabularx}

\subsubsection{Return Value}
0 on success, negative error code otherwise (\texttt{-ESRCH}, \texttt{-EINVAL}, \texttt{-EPERM}).

\subsubsection{Design Rationale}
Separating flag configuration from creation allows the caller to adjust permissions based on runtime conditions. This is more flexible than fixing flags at creation time.

\subsubsection{Changelog}
Introduced in 10.1.

\subsection{proc\_start}
\begin{lstlisting}
int proc_start(pid_t pid, const char *path, const parg_t *parg);
\end{lstlisting}

\subsubsection{Description}
Starts a process created with \texttt{proc\_new}. Loads the executable at \texttt{path} and passes \texttt{argc}/\texttt{argv}/\texttt{envp}. The process must not have been started before. Returns 0 on success, negative error code otherwise.

\subsubsection{Parameters}
\begin{tabularx}{\textwidth}{|l|l|X|}
\hline
\verb|pid| & \verb|pid_t| & PID of the process. \\
\verb|path| & \verb|const char *| & Path to the executable file. \\
\verb|parg| & \verb|const parg_t *| & Arguments and environment. \\
\hline
\end{tabularx}

\subsubsection{Return Value}
0 on success, negative error code (\texttt{-ENOENT}, \texttt{-ENOEXEC}, \texttt{-EINVAL}).

\subsection{proc\_terminate}
\begin{lstlisting}
_Noreturn void proc_terminate(void);
\end{lstlisting}
Terminates the calling process immediately. Does not return.

\subsection{proc\_kill}
\begin{lstlisting}
int proc_kill(pid_t pid);
\end{lstlisting}
Sends \texttt{PROCSIG\_KILL} to the process. Requires appropriate privilege. Returns 0 on success, negative error code.

\subsection{proc\_signal\_send}
\begin{lstlisting}
int proc_signal_send(pid_t pid, uint32_t signum);
\end{lstlisting}
Sends a signal (e.g., \texttt{PROCSIG\_TERM}) to a process. Returns 0 on success, negative error code.

\subsection{Orphan Process Handling}
When a process terminates, all its children are also terminated (they are not reparented to init). To keep a child alive after parent exits, the parent must have installed a signal handler for \texttt{PROCSIG\_CHLD} and the child must have been marked via an implementation‑defined mechanism. Default: children are killed.

\subsection{proc\_getpid, proc\_getppid}
\begin{lstlisting}
pid_t proc_getpid(void);
pid_t proc_getppid(void);
\end{lstlisting}
Return PID and parent PID of the caller.

\subsection{waitpid}
\begin{lstlisting}
pid_t waitpid(pid_t pid, int *status, int options);
\end{lstlisting}
Only \texttt{WNOHANG} supported. Returns child PID, 0 if \texttt{WNOHANG} and none exited, or \texttt{-ECHILD}.

\subsection{whoami}
\begin{lstlisting}
int whoami(void);
\end{lstlisting}
Returns the current ULA level (numeric).

\subsection{set\_privilege}
\begin{lstlisting}
int set_privilege(uint32_t mask);
\end{lstlisting}
Requires \texttt{hlud} authorization. Sets process privilege mask.

\section{Memory Management}

\subsection{page\_malloc}
\begin{lstlisting}
void *page_malloc(uint32_t count);
\end{lstlisting}

\subsubsection{Description}
Allocates \texttt{count} contiguous physical pages (4096 bytes each), zero-filled. Returns virtual address or \texttt{NULL}. Requires \texttt{PROC\_MEMORY} flag.

\subsubsection{Parameters}
\begin{tabularx}{\textwidth}{|l|l|X|}
\hline
\verb|count| & \verb|uint32_t| & Number of pages (must be > 0). \\
\hline
\end{tabularx}

\subsubsection{Return Value}
Success: non-NULL virtual address. Failure: \texttt{NULL}.

\subsubsection{Design Rationale}
Lowest-level memory primitive for DMA and page allocators. Zero-fill prevents information leaks.

\subsubsection{Changelog}
1.0 introduced; 5.0 added \texttt{PROC\_MEMORY} requirement; 10.2 renamed to \texttt{page\_malloc}/\texttt{page\_mfree}.

\subsection{page\_mfree}
\begin{lstlisting}
void page_mfree(void *pages, uint32_t count);
\end{lstlisting}
Frees pages allocated with \texttt{page\_malloc}. Behaviour undefined if parameters mismatch.

\subsection{shm\_open}
\begin{lstlisting}
fid_t shm_open(const char *name, int oflags, mode_t mode);
\end{lstlisting}
Creates/opens a shared memory object. Returns \texttt{fid\_t} or negative error.

\subsection{shm\_unlink}
\begin{lstlisting}
int shm_unlink(const char *name);
\end{lstlisting}
Removes a shared memory object. Returns 0 on success, negative error.
\section{File and Device I/O}

\subsection{open}
\begin{lstlisting}
fid_t open(const char *type, const char *mode, const char *addr);
\end{lstlisting}

\subsubsection{Description}
Opens a regular file (\texttt{type = "f"}) or a device (\texttt{type = "d"}). For network I/O, see \texttt{net\_*} functions (Section~\ref{sec:network}). Returns a file descriptor (\texttt{fid\_t}) that can be used with \texttt{read}, \texttt{write}, \texttt{io\_seek}, and \texttt{close}.

\subsubsection{Parameters}
\begin{tabularx}{\textwidth}{|l|l|X|}
\hline
\textbf{Parameter} & \textbf{Type} & \textbf{Description} \\
\hline
\verb|type| & \verb|const char *| & \texttt{"f"} (file) or \texttt{"d"} (device). \\
\verb|mode| & \verb|const char *| & \texttt{"r"} (read only), \texttt{"w"} (write only, truncate/create), \texttt{"rw"} (read/write). \\
\verb|addr| & \verb|const char *| & Path (absolute or relative) or device name (e.g., \texttt{/dev/tty0} or \texttt{"tty0"} interpreted under \texttt{/dev}). \\
\hline
\end{tabularx}

\subsubsection{Return Value}
Success: non-negative \texttt{fid\_t}. Failure: negative error code (\texttt{-ENOENT}, \texttt{-EPERM}, \texttt{-EINVAL}, \texttt{-EMFILE}).

\subsubsection{Design Rationale}
Uniform handling of files and devices simplifies the API. Network I/O is kept separate because its semantics differ significantly (no offsets, accept operation). The \texttt{type} string is used for type safety in the C library wrapper.

\subsubsection{Notes for Callers}
\begin{itemize}
    \item Requires \texttt{PROC\_IO\_READ} for reading, \texttt{PROC\_IO\_WRITE} for writing.
    \item \texttt{"w"} mode truncates the file to zero length. Use \texttt{"rw"} to preserve content.
\end{itemize}

\subsubsection{Changelog}
\begin{tabularx}{\textwidth}{|l|X|}
\hline
1.0 & Introduced. \\
5.0 & Added ULA permission checks. \\
6.0 & Removed network type \texttt{"n"} (moved to \texttt{net\_connect}). \\
10.1 & Added complete documentation. \\
\hline
\end{tabularx}

\subsection{close}
\begin{lstlisting}
int close(const char *type, fid_t fid);
\end{lstlisting}
Closes the I/O object. \texttt{type} must match.

\subsection{read}
\begin{lstlisting}
ssize_t read(const char *type, fid_t fid, void *buf, size_t count);
\end{lstlisting}
Reads up to \texttt{count} bytes from file/device/pipe. Returns bytes read, 0 on EOF, or negative error.

\subsection{write}
\begin{lstlisting}
ssize_t write(const char *type, fid_t fid, const void *buf, size_t count);
\end{lstlisting}
Writes up to \texttt{count} bytes. Returns bytes written or negative error.

\subsection{io\_seek}
\begin{lstlisting}
off_t io_seek(const char *type, fid_t fid, off_t offset, int whence);
\end{lstlisting}
Repositions file offset. \texttt{whence} = \texttt{SEEK\_SET}, \texttt{SEEK\_CUR}, \texttt{SEEK\_END}. Returns new offset or negative error.

\subsection{godir, godir\_byfid, getwd}
Working directory functions as in version 9.0.

\subsection{pipe}
\begin{lstlisting}
fid_t pipe(pid_t to_who);
\end{lstlisting}
Creates a unidirectional pipe. Write end returned; read end automatically transferred to \texttt{to\_who}. Type is \texttt{"p"}.

\subsection{cpfid, cpfid2, set\_flags, umask}
Standard descriptor duplication and flag functions.

\section{Network Communications}
\label{sec:network}

Network I/O is handled separately from file I/O. All functions require the \texttt{PROC\_NETWORK} flag.

\subsection{net\_connect}
\begin{lstlisting}
fid_t net_connect(const char *protocol, const char *address);
\end{lstlisting}

\subsubsection{Description}
Establishes a connection to a remote host. \texttt{protocol} is \texttt{"tcp"} (future: \texttt{"udp"}). \texttt{address} is \texttt{"ip:port"} (e.g., \texttt{"192.168.1.1:80"}).

\subsubsection{Parameters}
\begin{tabularx}{\textwidth}{|l|l|X|}
\hline
\verb|protocol| & \verb|const char *| & \texttt{"tcp"} (or \texttt{"udp"} in future). \\
\verb|address| & \verb|const char *| & e.g., \texttt{"192.168.1.1:80"} or \texttt{"[::1]:8080"}. \\
\hline
\end{tabularx}

\subsubsection{Return Value}
Success: \texttt{fid\_t}. Failure: negative error (\texttt{-ECONNREFUSED}, \texttt{-EINVAL}, \texttt{-ENETUNREACH}).

\subsubsection{Design Rationale}
Separate from file I/O to avoid confusion (no file offsets, accept semantics). Cleaner API for network-specific operations.

\subsubsection{Changelog}
6.0 introduced; 10.1 added complete documentation.

\subsection{net\_listen}
\begin{lstlisting}
fid_t net_listen(const char *protocol, const char *address);
\end{lstlisting}
Creates a listening socket. Returns \texttt{fid\_t} or negative error (\texttt{-EADDRINUSE}, \texttt{-EINVAL}).

\subsection{net\_accept}
\begin{lstlisting}
fid_t net_accept(fid_t listen_fid);
\end{lstlisting}
Accepts an incoming connection. Returns new \texttt{fid\_t} or negative error.

\subsection{net\_send}
\begin{lstlisting}
ssize_t net_send(fid_t fid, const void *buf, size_t count);
\end{lstlisting}
Sends data on a connected socket. Returns bytes sent or negative error (\texttt{-EPIPE}, \texttt{-ENOTCONN}).

\subsection{net\_recv}
\begin{lstlisting}
ssize_t net_recv(fid_t fid, void *buf, size_t count);
\end{lstlisting}
Receives data from a connected socket. Returns bytes received, 0 on remote close, or negative error.

\subsection{net\_disconnect}
\begin{lstlisting}
int net_disconnect(fid_t fid);
\end{lstlisting}
Closes the network connection. Returns 0 on success, negative error.

\section{Event Multiplexing (hear Interface)}
\label{sec:hear}

The \texttt{hear} interface allows waiting for events from multiple sources (I/O, IPC, signals, process state).

\subsection{hear\_add}
\begin{lstlisting}
int hear_add(fid_t fid, int events);
\end{lstlisting}

\subsubsection{Description}
Adds \texttt{fid} to the monitored set. \texttt{events} is a bitmask of \texttt{HEAR\_IN}, \texttt{HEAR\_OUT}, \texttt{HEAR\_IPC}, \texttt{HEAR\_SIGNAL}, \texttt{HEAR\_PROC}, etc.

\subsubsection{Parameters}
\begin{tabularx}{\textwidth}{|l|l|X|}
\hline
\verb|fid| & \verb|fid_t| & File/pipe/socket/process to monitor. \\
\verb|events| & \verb|int| & Bitmask of event types. \\
\hline
\end{tabularx}

\subsubsection{Return Value}
0 on success, negative error otherwise.

\subsubsection{Design Rationale}
Unified event notification simplifies application design. One interface for all async events.

\subsubsection{Changelog}
6.0 introduced; 10.1 added complete documentation.

\subsection{hear\_remove}
\begin{lstlisting}
int hear_remove(fid_t fid);
\end{lstlisting}
Removes \texttt{fid} from the monitored set.

\subsection{hear\_filter}
\begin{lstlisting}
int hear_filter(fid_t fid, int new_events);
\end{lstlisting}
Changes the event mask for an existing source.

\subsection{hear\_wait}
\begin{lstlisting}
int hear_wait(void *events_out, int max_events, int timeout_ms);
\end{lstlisting}

\subsubsection{Description}
Waits for events among registered sources. Fills \texttt{events\_out} with up to \texttt{max\_events} event structures. Each event contains \texttt{handle}, \texttt{type}, and \texttt{userdata} (if set).

\subsubsection{Parameters}
\begin{tabularx}{\textwidth}{|l|l|X|}
\hline
\verb|events_out| & \verb|void *| & Pointer to array of \texttt{struct hear\_event}. \\
\verb|max_events| & \verb|int| & Maximum number of events to return. \\
\verb|timeout_ms| & \verb|int| & Timeout in milliseconds (0 = non-blocking, -1 = infinite). \\
\hline
\end{tabularx}

\subsubsection{Return Value}
Number of events received, 0 on timeout, negative error on failure.

\subsubsection{Example}
\begin{lstlisting}
fid_t fd1 = open("f", "r", "/tmp/a");
fid_t fd2 = open("f", "r", "/tmp/b");
hear_add(fd1, HEAR_IN, NULL);
hear_add(fd2, HEAR_IN, NULL);
struct hear_event ev[4];
int n = hear_wait(ev, 4, 1000);
for (int i = 0; i < n; i++) {
    if (ev[i].handle == fd1) read("f", fd1, ...);
    else if (ev[i].handle == fd2) read("f", fd2, ...);
}
\end{lstlisting}
\section{Inter-Process Communication (IPC)}

SLFD provides synchronous message passing between threads. Each thread has a kernel‑managed message queue.

\subsection{System Calls}
\begin{lstlisting}
int sys_ipc_send(thread_id target, ipc_message_t *msg);
int sys_ipc_receive(ipc_message_t *msg);
\end{lstlisting}

\subsubsection{Description}
\texttt{sys\_ipc\_send} sends \texttt{msg} to the target thread. The kernel fills \texttt{source} and \texttt{timestamp}. If the target's queue is full, the caller may block (implementation‑defined). \texttt{sys\_ipc\_receive} retrieves the next message from the caller's queue. If no message is available, the call blocks until one arrives.

\subsubsection{Return Value}
0 on success, negative error code otherwise.

\subsection{C Library Wrappers}
\begin{lstlisting}
int ipc_send_msg(uint32_t type, uint32_t flags, uint32_t receiver_pid,
                 uint32_t data_len, const void *data);
int ipc_receive_msg(ipc_message_t *msg, int timeout_ms);
\end{lstlisting}
Wrapper functions for convenience. \texttt{ipc\_send\_msg} sends to the main thread of the target process. \texttt{ipc\_receive\_msg} supports timeout (0 = non‑blocking, -1 = infinite).

\subsection{Error Codes for IPC}
\texttt{-EAGAIN}, \texttt{-ETIMEDOUT}, \texttt{-EMSGSIZE}, \texttt{-ENOMSG}.

\subsection{IPC Performance Tuning (Implementation‑Defined)}
\begin{lstlisting}
int ipc_set_buffer(thread_id target, void *buffer, size_t size);
int ipc_use_buffer(thread_id target, void *buffer, size_t size);
\end{lstlisting}
Optional functions for shared memory message caching. May be no‑ops.

\section{Thread Management}
\label{sec:threads}

\subsection{thread\_create}
\begin{lstlisting}
int thread_create(thread_id *tid, void *(*start)(void*), void *arg);
\end{lstlisting}
Creates a new thread that starts executing \texttt{start(arg)}. Inherits signal mask and ULA privilege. Returns 0 on success.

\subsection{thread\_join}
\begin{lstlisting}
int thread_join(thread_id tid, void **retval);
\end{lstlisting}
Waits for thread termination. Stores exit value in \texttt{retval}.

\subsection{thread\_exit}
\begin{lstlisting}
_Noreturn void thread_exit(void *retval);
\end{lstlisting}
Terminates the calling thread.

\subsection{Mutexes and Condition Variables}
Standard POSIX‑style primitives:
\begin{lstlisting}
typedef struct mutex { /* opaque */ } mutex_t;
int mutex_init(mutex_t *mutex);
int mutex_lock(mutex_t *mutex);
int mutex_unlock(mutex_t *mutex);
int mutex_destroy(mutex_t *mutex);

typedef struct cond { /* opaque */ } cond_t;
int cond_init(cond_t *cond);
int cond_wait(cond_t *cond, mutex_t *mutex);
int cond_signal(cond_t *cond);
int cond_broadcast(cond_t *cond);
int cond_destroy(cond_t *cond);
\end{lstlisting}
All return 0 on success, negative error otherwise.

\section{Error Codes}
\label{sec:errors}

Full list of error codes (as in version 9.0). Include all codes with descriptions. I'll provide the complete table (longtable).

\begin{longtable}{|l|l|p{0.55\textwidth}|}
\hline
\textbf{Symbolic Name} & \textbf{Value} & \textbf{Description} \\
\hline
\endfirsthead
\hline
\multicolumn{3}{|c|}{\textit{Continued from previous page}} \\
\hline
\textbf{Symbolic Name} & \textbf{Value} & \textbf{Description} \\
\hline
\endhead
\hline
\multicolumn{3}{|r|}{\textit{Continued on next page}} \\
\endfoot
\endlastfoot

\hline
\verb|-EPERM| & 1 & Permission denied. \\
\verb|-ENOENT| & 2 & No such file or directory. \\
\verb|-ESRCH| & 3 & No such process. \\
\verb|-EINTR| & 4 & Interrupted system call. \\
\verb|-EIO| & 5 & I/O error. \\
\verb|-ENXIO| & 6 & No such device or address. \\
\verb|-E2BIG| & 7 & Argument list too long. \\
\verb|-ENOEXEC| & 8 & Exec format error. \\
\verb|-EBADF| & 9 & Bad file descriptor. \\
\verb|-ECHILD| & 10 & No child process. \\
\verb|-EAGAIN| & 11 & Resource temporarily unavailable. \\
\verb|-ENOMEM| & 12 & Out of memory. \\
\verb|-EACCES| & 13 & Permission denied (file access). \\
\verb|-EFAULT| & 14 & Bad address. \\
\verb|-ENOTBLK| & 15 & Block device required. \\
\verb|-EBUSY| & 16 & Device or resource busy. \\
\verb|-EEXIST| & 17 & File exists. \\
\verb|-EXDEV| & 18 & Cross-device link. \\
\verb|-ENODEV| & 19 & No such device. \\
\verb|-ENOTDIR| & 20 & Not a directory. \\
\verb|-EISDIR| & 21 & Is a directory. \\
\verb|-EINVAL| & 22 & Invalid argument. \\
\verb|-ENFILE| & 23 & File table overflow. \\
\verb|-EMFILE| & 24 & Too many open files. \\
\verb|-ENOTTY| & 25 & Inappropriate I/O control. \\
\verb|-ETXTBSY| & 26 & Text file busy. \\
\verb|-EFBIG| & 27 & File too large. \\
\verb|-ENOSPC| & 28 & No space left on device. \\
\verb|-ESPIPE| & 29 & Illegal seek. \\
\verb|-EROFS| & 30 & Read-only file system. \\
\verb|-EMLINK| & 31 & Too many links. \\
\verb|-EPIPE| & 32 & Broken pipe. \\
\verb|-EDOM| & 33 & Math argument out of domain. \\
\verb|-ERANGE| & 34 & Result too large. \\
\verb|-EDEADLK| & 35 & Deadlock avoided. \\
\verb|-ENAMETOOLONG| & 36 & File name too long. \\
\verb|-ENOLCK| & 37 & No locks available. \\
\verb|-ENOSYS| & 38 & Function not implemented. \\
\verb|-ENOTEMPTY| & 39 & Directory not empty. \\
\verb|-ELOOP| & 40 & Too many symbolic links. \\
\verb|-EWOULDBLOCK| & 41 & Same as -EAGAIN. \\
\verb|-ENOMSG| & 42 & No message of desired type. \\
\verb|-EIDRM| & 43 & Identifier removed. \\
\verb|-ETIMEDOUT| & 44 & Connection timed out. \\
\verb|-ERESTART| & 45 & Restart syscall (internal). \\
\verb|-EMSGSIZE| & 46 & Message too long. \\
\verb|-EADDRINUSE| & 47 & Address already in use. \\
\verb|-ECONNREFUSED| & 48 & Connection refused. \\
\verb|-ECONNRESET| & 49 & Connection reset by peer. \\
\verb|-ENOTCONN| & 50 & Socket not connected. \\
\verb|-ETIME| & 51 & Timer expired. \\
\verb|-ECANCELED| & 52 & Operation canceled. \\
\hline
\end{longtable}

\appendix

\section{ULA: User Level Assessment}
\label{app:ula}

This appendix defines the ULA privilege model. It replaces traditional Unix UID/GID and mode\_t with a level-based bitmask system.

\subsection{Privilege Levels}
\begin{lstlisting}
#define LV_HIGHEST   0
#define LV_GENERAL   1
#define LV_LOW       2
#define LV_GETTING   3
\end{lstlisting}

\subsection{File Attribute Mask Layout}
The 32-bit mask:
\begin{itemize}
    \item Bits 31–24: creator level
    \item Bits 23–16: caller level
    \item Bits 15–8: flags (EXEC, HIDDEN, ARCHIVE)
    \item Bits 7–4: file type (FILE, DIR, LINK, DEV)
    \item Bits 3–0: reserved
\end{itemize}

\subsection{Access Control Matrix}
\begin{table}[htbp]
\centering
\small
\begin{tabularx}{\textwidth}{|l|X|c|c|c|}
\hline
\textbf{Operation} & \textbf{LV\_HIGHEST} & \textbf{LV\_GENERAL} & \textbf{LV\_LOW} & \textbf{LV\_GETTING} \\
\hline
\texttt{proc\_kill}(any process) & Yes & No & No & No \\
\texttt{proc\_kill}(own process) & Yes & Yes & Yes & Yes \\
\texttt{proc\_new("s", ...)} (clone) & Yes & Yes & No & No \\
\texttt{proc\_signal\_send}(any) & Yes & No & No & No \\
\texttt{open("d", ...)} (any device) & Yes & Limited\textsuperscript{a} & No & Limited\textsuperscript{b} \\
\texttt{open("f", "w")} on system dirs & Yes & No & No & No \\
\texttt{open("f", "r")} on system files & Yes & No\textsuperscript{c} & No & Yes\textsuperscript{d} \\
\texttt{proc\_new("n", ...)} & Yes & Yes & Yes & Yes \\
\texttt{getwd}, \texttt{stat}, \texttt{mkdir} & Yes & Own files only & Own files only (restricted) & Own work dir\textsuperscript{e} \\
\hline
\end{tabularx}
\end{table}
Notes:
\begin{enumerate}
    \item \texttt{LV\_GENERAL} may open non‑sensitive devices (e.g., \texttt{/dev/tty}) but not raw block devices.
    \item \texttt{LV\_GETTING} may open devices required for its service (e.g., \texttt{/dev/null}) but not console.
    \item \texttt{LV\_GENERAL} cannot read system config files (e.g., \texttt{/conf/shadow}).
    \item \texttt{LV\_GETTING} may read specific system files as defined by service manifest.
    \item \texttt{LV\_GETTING} uses service‑specific working directory (e.g., \texttt{/usr/myapp}) and cannot access others' home.
\end{enumerate}

\subsection{Textual Database Example}
\begin{lstlisting}
# /conf/users
alice:{staff-lvgeneral}+{audio-lvlow}
bob:{staff-lvlow}
root:{wheel-lvhighest}
# /conf/groups
staff:500:alice,bob
wheel:0:root
audio:300:alice
\end{lstlisting}

\subsection{Changelog}
\begin{tabularx}{\textwidth}{|l|X|}
\hline
1.0 & Initial design. \\
3.0 & Added access control matrix. \\
5.0 & Bitmask layout introduced. \\
7.0 & Added textual database example. \\
9.0 & Added rationale. \\
10.1 & Moved to Appendix A. \\
\hline
\end{tabularx}

\section{SLFD Free Desktop (SLFD FD)}
\label{app:sfd}

This appendix defines the root directory layout for SLFD‑based systems.

\subsection{Root Directory Structure}
\begin{description}
    \item[\texttt{/bin}] Essential system binaries (e.g., \texttt{sh}, \texttt{ls}).
    \item[\texttt{/dev}] Device special files.
    \item[\texttt{/conf}] System and application configuration files (replaces \texttt{/etc}).
    \item[\texttt{/usr}] Combined user and system resources (user home directories + system binaries/libraries).
    \item[\texttt{/lib}] Common libraries (may be symlink to \texttt{/usr/lib}).
    \item[\texttt{/tmp}] Temporary files (cleared on reboot).
    \item[\texttt{/sys}] System files (e.g., \texttt{/sys/kernel/debug}).
    \item[\texttt{/mnt}] Mount point for devices.
\end{description}

\subsection{Rationale}
Avoids FHS and freedesktop.org complexity. All configuration in \texttt{/conf}, all data under \texttt{/usr}, essential binaries in \texttt{/bin}.

\subsection{Example}
\begin{lstlisting}
/
+-- bin/
|   +-- sh
|   +-- ls
+-- dev/
|   +-- tty0
|   +-- random
|   +-- null
+-- conf/
|   +-- passwd
|   +-- shadow
|   +-- slfd.conf
|   +-- wahlu
+-- usr/
|   +-- bin/    (optional symlink to /bin)
|   +-- lib/
|   +-- alice/
|   +-- bob/
|   +-- ...
+-- mnt/
|   +-- dev1/
|   +-- dev2/
|   +-- ...
+-- sys/
|   +-- ...
+-- lib/        (optional symlink to usr/lib)
+-- tmp/
\end{lstlisting}

\subsection{Changelog}
\begin{tabularx}{\textwidth}{|l|X|}
\hline
1.0 & Initial: \texttt{/bin}, \texttt{/dev}, \texttt{/etc}, \texttt{/usr}, \texttt{/lib}, \texttt{/tmp}. \\
2.0 & Replaced \texttt{/etc} with \texttt{/conf}. \\
3.0 & Added \texttt{/sys}, \texttt{/mnt}. \\
5.0 & Renamed to SLFD FD. \\
10.1 & Moved to Appendix B. \\
\hline
\end{tabularx}

\section{SLFD Shell Standard (3S)}
\label{app:3s}

Non‑mandatory shell standard. Commands:
\begin{longtable}{|l|p{0.7\textwidth}|}
\hline
\textbf{Command} & \textbf{Description} \\
\hline
\texttt{lookup [path]} & List directory contents. \\
\texttt{goto <path>} & Change directory. \\
\texttt{where} & Print current directory. \\
\texttt{create.dir <path>} & Create directory. \\
\texttt{del.dir <path>} & Remove empty directory. \\
\texttt{del <path>} & Remove file/link. \\
\texttt{connect <file...>} & Concatenate files to stdout. \\
\texttt{count [options] [file]} & Count lines/words/bytes. \\
\texttt{head [-n N] [file]} & Output first N lines. \\
\texttt{tail [-n N] [file]} & Output last N lines. \\
\texttt{filter <pattern> [file...]} & Search for pattern. \\
\texttt{pm [options]} & Process management. \\
\texttt{pm kill <PID> [signal]} & Send signal. \\
\texttt{env [var[=value]...]} & Environment. \\
\texttt{sysinfo} & System information. \\
\texttt{info-of <command>} & Manual page. \\
\hline
\end{longtable}

Extended ULA commands: \texttt{ula file-change}, \texttt{ula user-change}.

\subsection{Non‑mandatory Status}
This standard is optional. A shell may claim 3S compliance if it implements these commands.

\subsection{Changelog}
\begin{tabularx}{\textwidth}{|l|X|}
\hline
1.0 & Initial. \\
4.0 & Updated to match SLFD FD naming. \\
5.0 & Added \texttt{pm}, \texttt{pm kill}. \\
10.1 & Moved to Appendix C. \\
\hline
\end{tabularx}

\section{Kernel Primitives and System Call Interfaces}
\label{app:syscalls}

This appendix lists mandatory and optional kernel primitives.

\subsection{Mandatory Kernel Primitives}
Following the MINIX 3 model, the kernel itself provides only the minimal
primitives needed to manipulate process control blocks (PCBs) and pass
messages. Higher-level process management policy (validation, sequencing,
signal delivery semantics) is implemented by a user-space process server;
see \S\ref{app:syscalls}.2.
\begin{longtable}{|l|p{0.65\textwidth}|}
\hline
\textbf{Primitive} & \textbf{Description} \\
\hline
\texttt{\detokenize{sys_pcb_alloc}} & Allocate a PCB slot and PID; does not load or start execution. Used internally by the process server. \\
\texttt{\detokenize{sys_pcb_free}} & Release a PCB slot. Used internally by the process server. \\
\texttt{sys\_mm\_alloc\_pages} & Allocate physical pages (backing \texttt{page\_malloc}). \\
\texttt{sys\_mm\_free\_pages} & Free pages (backing \texttt{page\_mfree}). \\
\texttt{sys\_shm\_open} & Open shared memory. \\
\texttt{sys\_shm\_unlink} & Remove shared memory. \\
\texttt{sys\_ipc\_send} & Send IPC message. \\
\texttt{sys\_ipc\_receive} & Receive IPC message. \\
\texttt{sys\_getrlimit} & Get resource limits. \\
\texttt{sys\_setrlimit} & Set resource limits. \\
\texttt{sys\_sleep\_pro} & Sleep for seconds. \\
\texttt{sys\_timer\_crate} & Create timer. \\
\texttt{sys\_get\_fucktimer} & Get monotonic time. \\
\texttt{sys\_settimeofday} & Set real‑time clock. \\
\texttt{sys\_info\_sys} & Fill \texttt{utsname}. \\
\texttt{sys\_getrandom} & Fill buffer with random. \\
\texttt{sys\_hlud\_request} & Request elevation. \\
\hline
\end{longtable}

\subsection{Optional Primitives}
These may be implemented as user‑space servers.
\begin{longtable}{|p{0.35\textwidth}|p{0.55\textwidth}|}
\hline
\textbf{Function} & \textbf{Typical Implementation} \\
\hline
\texttt{proc\_new}, \texttt{proc\_setflags}, \texttt{proc\_start}, \texttt{proc\_terminate}, \texttt{proc\_kill}, \texttt{proc\_signal\_send}, \texttt{proc\_getpid}, \texttt{proc\_getppid}, \texttt{waitpid}, \texttt{whoami}, \texttt{set\_privilege} & Process management server (PM‑equivalent, per the MINIX 3 model)... \\
\texttt{proc\_new}, \texttt{proc\_setflags}, \texttt{proc\_start}, \texttt{proc\_terminate}, \texttt{proc\_kill}, \texttt{proc\_signal\_send}, \texttt{proc\_getpid}, \texttt{proc\_getppid}, \texttt{waitpid}, \texttt{whoami}, \texttt{set\_privilege} & Process management server (PM‑equivalent, per the MINIX 3 model). Implements all argument validation, ULA privilege checks, and multi‑step sequencing; internally uses \texttt{sys\_pcb\_alloc}/\texttt{sys\_pcb\_free} for the actual kernel-level PCB manipulation. Clients reach this server over IPC, never via direct trap. \\
\texttt{\detokenize{open, close, read, write, io_seek}} & File system server. \\
\texttt{\detokenize{godir, godir_byfid, getwd}} & File system server. \\
\texttt{pipe} & IPC or kernel service. \\
\texttt{cpfid}, \texttt{cpfid2} & Kernel or server. \\
\texttt{set\_flags} & Kernel. \\
\texttt{umask} & User‑space library. \\
\texttt{dev\_control} & Device driver server. \\
\texttt{chmod}, \texttt{chown}, \texttt{stat} & File system server. \\
\texttt{mkdir}, \texttt{rmdir} & File system server. \\
\texttt{symlink}, \texttt{readlink} & File system server. \\
\texttt{mount}, \texttt{umount} & File system server. \\
\texttt{reboot} & System management server. \\
\texttt{fs\_*} & File system server. \\
\texttt{thread\_*}, mutex/cond & Library or kernel-assisted. \\
\texttt{cp0\_*}, \texttt{aio\_*} & Implementation-defined. \\
\texttt{\detokenize{ipc_set_buffer}}, \texttt{\detokenize{ipc_use_buffer}} & Implementation-defined. \\
\texttt{hear\_*} & System call or event server. \\
\texttt{net\_*} & Network server. \\
\hline
\end{longtable}

\subsection{Changelog}
\begin{tabularx}{\textwidth}{|l|X|}
\hline
1.0 & Initial. \\
3.0 & Added \texttt{sys\_proc\_signal\_send}. \\
5.0 & Added ULA primitives. \\
6.0 & Added hear and network. \\
7.0 & Split into mandatory/optional. \\
10.1 & Corrected names, moved to Appendix D. \\
10.2 & Adopted MINIX 3 model: replaced monolithic \texttt{sys\_proc\_*} kernel primitives with minimal \texttt{sys\_pcb\_alloc}/\texttt{sys\_pcb\_free}; process management (\texttt{proc\_*}, \texttt{waitpid}, \texttt{whoami}, \texttt{set\_privilege}) moved to Optional Primitives as a user‑space process server. \\
\hline
\end{tabularx}

\end{document}
